% =============================================================================
% Complexity vs. efficiency analysis
% Drafted per anthropic-skills:polishing-manuscripts. Every number traces to
% analysis/out/inference_time_per_model.csv (parameters, inference latency),
% analysis/out/flops_per_model.csv (FLOPs, via analysis/flop_utils.py), or
% the already-verified CMAPSS results tables (RMSE, Score).
% Evidence-vs-interpretation audit at the end of this file.
% =============================================================================

\subsection{Complexity vs. efficiency analysis}
\label{subsec:complexity}

Table~\ref{tab:complexity} and Figure~\ref{fig:complexity_rmse}/\ref{fig:complexity_score}
summarize each method's parameter count and FLOPs against its mean CMAPSS
RMSE and Score (all 12 source$\to$target pairs, all seeds), with inference
latency measured directly on the trained checkpoints -- one full test-set
forward pass per timed run, averaged over repeated runs, entirely separate
from training cost -- rather than taken from a training-time proxy, since
the two are not interchangeable and only the former reflects deployment
cost. FLOPs were computed on a single batch-size-1 sample per model via
\texttt{thop}, extended with a monkey-patched einsum/matmul/FFT counter for
the custom graph and spectral operations \texttt{thop}'s standard layer
hooks cannot see (Section~\ref{subsec:flops_method}); reported values are
MACs$\times2$.

\begin{table}[t]
\centering
\caption{Model complexity and CMAPSS efficiency. Parameters and FLOPs for
LBGN are pair-weighted means across its two source-dependent architecture
variants (Section~\ref{subsec:flops_method}); every other model's
architecture, and therefore its parameter count and FLOPs, is fixed across
all 12 pairs. \textbf{Bold} marks the best (lowest) value per column.}
\label{tab:complexity}
\begin{tabular}{lrrrrr}
\toprule
Model & Parameters & FLOPs (M) & Inference (ms) & RMSE & Score \\
\midrule
LBGN (ours)              & \textbf{11{,}538}    & \textbf{0.47} & 8.00           & \textbf{18.15} & \textbf{2{,}007.75} \\
DAGCN~\cite{li2021DAGCN}       & 1{,}586{,}447   & 4.32          & 9.61           & 26.61          & 31{,}993.92 \\
EviAdapt~\cite{hou2025EviAdapt} & 121{,}480       & 6.93          & 11.35          & 25.86          & 6{,}093.46 \\
DAST~\cite{li2022DAST}         & 151{,}652       & 9.31          & 6.32           & 25.74          & 4{,}496.76 \\
CADA~\cite{ragab2021CADA}      & 115{,}265       & 6.92          & 11.29          & 25.95          & 5{,}467.99 \\
TACDA~\cite{hou2025TACDA}      & 124{,}505       & 7.51          & 22.17          & 26.61          & 6{,}546.00 \\
OCS-DANN~\cite{nejjar2023OCSDANN} & 26{,}405     & 2.33          & \textbf{1.43}  & 25.38          & 9{,}935.86 \\
MDAN~\cite{furqon2024MDAN}     & 65{,}089        & 6.41          & 15.80          & 25.73          & 5{,}932.85 \\
CCDG~\cite{mohamed2022CCDG}    & 353{,}921       & 4.44          & 1.94           & 31.12          & 22{,}760.47 \\
\bottomrule
\end{tabular}
\end{table}

LBGN uses 11{,}538 parameters and 0.47~MFLOPs on average, 2.3--137.5$\times$
fewer parameters and 4.9--19.7$\times$ fewer FLOPs than every compared
baseline (Table~\ref{tab:complexity}); OCS-DANN is next-smallest on both
measures (26{,}405 parameters, 2.33~MFLOPs). Neither figure is perfectly
fixed across source domains: LBGN's LCE\_dim is tuned per source
(Section~\ref{subsec:flops_method}), giving 13{,}716 parameters / 0.52~MFLOPs
on FD001 and 10{,}812 / 0.45~MFLOPs elsewhere, so the values reported are
pair-weighted means (3 of 12 pairs use the FD001 variant) -- consistent with
how RMSE, Score, and inference time are already averaged over all 12 pairs.
Despite this reduction, LBGN attains the lowest mean RMSE (18.15) and the
lowest mean Score (2{,}007.75) of all nine compared methods
(Section~\ref{subsec:results}), so the accuracy gain does not come from
added model capacity.

Parameters and FLOPs do not rank the same way across baselines, which
matters for what "complexity" is taken to mean. DAGCN has by far the most
parameters (1{,}586{,}447, 137.5$\times$ LBGN) but ranks only third-lowest on
FLOPs (4.32~MFLOPs) -- lower than DAST, TACDA, EviAdapt, CADA, and MDAN,
each of which has 10--24$\times$ fewer parameters (10.5$\times$ for DAST,
24.4$\times$ for MDAN). Conversely, DAST has the
third-largest parameter count (151{,}652) -- an order of magnitude below
DAGCN and well below CCDG, but still more than every other baseline -- yet
the highest FLOPs of any compared method (9.31~MFLOPs). This is consistent with
DAGCN's cost sitting mostly in wide layers that are evaluated once per
forward pass, while DAST's and TACDA's attention-based components scale
with sequence length rather than channel width, so a parameter-count
comparison alone would rank these methods' computational cost in the wrong
order; FLOPs and parameters are both necessary here; the table reports both
rather than treating one as a proxy for the other.

On inference latency, LBGN is competitive but not the fastest: it ranks
fourth of nine methods (8.00~ms/run), behind OCS-DANN (1.43~ms), CCDG
(1.94~ms), and DAST (6.32~ms), and ahead of DAGCN, CADA, EviAdapt, MDAN, and
TACDA (9.61--22.17~ms). The two faster methods do not dominate LBGN
overall, however: OCS-DANN's RMSE (25.38) and Score (9{,}935.86) and CCDG's
RMSE (31.12) and Score (22{,}760.47) are both substantially worse than
LBGN's, so their latency advantage -- on the order of a few milliseconds in
absolute terms -- trades against a large accuracy gap rather than a
comparable one. We report LBGN's latency rank plainly rather than
characterizing it as fastest, since it is not; the evidence instead supports
a parameter- and FLOPs-efficiency claim specifically (best accuracy at the
lowest parameter count and lowest FLOPs of any compared method), not a
latency-efficiency claim.

% =============================================================================
% AUDIT NOTE (remove before submission; not part of the manuscript text)
%
% Directly measured (analysis/out/inference_time_per_model.csv, generated by
% scripts/measure_inference_time_cmapss.py -- real forward-pass-only timing,
% NOT the CSV's mislabeled 'inference_time' column, which is actually mean
% training-epoch time):
%   - params_mean and mean_ms for all 9 compared models
%   - the 1st-9th latency ranking and the "LBGN ranks 4th" claim
%   - LBGN's 10,812-13,716 parameter range (configs/hparams.py's per-source
%     LCE_dim)
%
% Directly measured (analysis/out/flops_per_model.csv, generated by
% scripts/measure_flops_cmapss.py + analysis/flop_utils.py -- thop for
% standard layers, monkey-patched einsum/matmul/FFT for the custom ops thop
% cannot see):
%   - total_macs_incl_fft x2 for all 9 models
%   - LBGN's two FLOPs variants (0.52M FD001 / 0.45M elsewhere), pair-weighted
%     to 0.47M for consistency with the params_mean weighting above -- NOTE
%     this replaces the simple 2-variant average (0.489M) given in an earlier
%     turn of this session; the pair-weighted figure is the one now used
%     throughout, both here and in the array-style data given previously
%     should be corrected to match if reused elsewhere.
%   - the parameters-vs-FLOPs rank divergence for DAGCN (1st on params,
%     3rd-lowest on FLOPs) and DAST (7th on params, 1st on FLOPs) -- computed
%     directly from the CSV, not estimated
%
% Directly measured (build_final_local_logs_table.py's verified
% load_all()/collect(), same numbers already in the CMAPSS results tables):
%   - all RMSE/Score means
%
% Interpretation (explicitly flagged as such above, not asserted as fact):
%   - "the accuracy gain does not come from added model capacity" -- true for
%     DAGCN specifically and consistent with the overall pattern, but not a
%     controlled ablation isolating parameter count as a causal variable.
%   - the mechanistic explanation for DAGCN-vs-DAST/TACDA's FLOPs/params
%     divergence ("wide layers evaluated once" vs. "attention scaling with
%     sequence length") is architecturally plausible given each model's
%     known design (Chebyshev graph conv for DAGCN; self-attention for
%     DAST/TACDA) but was not verified by a layer-by-layer FLOPs breakdown --
%     flag as informed inference, not a measured decomposition, unless such
%     a breakdown is added.
%   - the "trades against a large accuracy gap" framing for OCS-DANN/CCDG is
%     a reading of the same table, not a separate measurement.
%
% Open dependencies to resolve before this compiles:
%   - \label{fig:complexity_rmse} / \label{fig:complexity_score} -- must be
%     set on the \includegraphics figure environments wrapping
%     figures/Inference_complexity_rmse.eps and
%     figures/Inference_complexity_score.eps. NOTE: those two figures'
%     x-axis is currently FLOPs (changed from inference time earlier this
%     session), so their captions should describe FLOPs, not latency.
%   - \label{subsec:results} -- points at wherever the CMAPSS RMSE/Score
%     results subsection lives; update if its actual label differs.
%   - \label{subsec:flops_method} -- points at wherever FLOPs/parameter
%     counting methodology is (or should be) described; this section
%     references it twice but does not itself define it. Add a methods
%     paragraph or point this at the right existing subsection.
% =============================================================================
